\begin{abstract}
We determine the surface obtained by Grothendieck's corrected disk gluing
for every finite non-pencil arrangement of distinct pseudolines in the
projective plane, allowing multiple intersections. We use the variant
$\mathcal X_2$ of cote 154, scans 5 and 7. If $c_j$ counts the ordinary
unoriented relative positions passing through exactly $j$ old vertices,
then
\[
 \mathcal X_2\cong
 \#^{\,n+1+4\sum_{k\ge1}k c_{2k+1}}\mathbb {RP}^{2}.
\]
We prove that the gluing is well defined, every closed cell embeds,
and every vertex link is a four-cycle. The extension quotient is
homotopy equivalent to $\mathbb {RP}^{2}$; its specified real blow-up
at the parallel extensions is homotopy equivalent to
$\#^{n+1}\mathbb {RP}^{2}$. The dual manuscript surface is an explicitly
described subcomplex of that blow-up, and the inclusion is a homotopy
equivalence exactly when no ordinary position passes through three old
vertices. A fully enumerated six-line arrangement has
$(s_0,s_1,s_2)=(360,720,315)$ and genus 47, refuting the proposed
universal Euler formula. We also give an exact finite image criterion
for deployed polygons, reconstruct the weighted cellular graph,
prove the disk sufficiency and cluster-reversal statements, and
separate the boundary-to-interior tethers of cote 156-4 from proper
chords. All finite counterexamples come with exact code and outputs.
\end{abstract}

\tableofcontents
\newpage
\section*{Scope, conventions, and outcome}
\addcontentsline{toc}{section}{Scope, conventions, and outcome}
We work with labelled, finite arrangements of distinct tame pseudolines
in $\mathbb {RP}^{2}$. Each pair meets at precisely one transverse
crossing. ``Nontrivial'' means that the arrangement is not a pencil,
so its oriented matroid has rank three. The word ``simple'' for the
oriented matroid excludes loops and parallel old elements; a
\emph{simple arrangement} additionally excludes triple intersections.
The latter extra condition is imposed only in the explicitly named
finite computations. Relative positions are combinatorial isotopy
classes, equivalently labelled extension data under the relative
representation comparison proved in A.3. We do not count parametrized
embedded curves or points in a metric realization space.

The construction in cote 154 has several versions. Our $\mathcal X$
always means $\mathcal X_2$: every ordinary relative position labels a
face, and exceptional sweeps across old pseudolines are included.
The capped variant $\mathcal X_0$ is different. A vertex of our primal
surface has valence four; correspondingly a face of the dual graph is
a quadrangle. The primal faces themselves have arbitrary polygonal
orders. This distinction resolves the potentially ambiguous phrase
``faces of order four'' in the supplied question.

\begin{center}\small
\begin{tabular}{p{0.16\linewidth}p{0.62\linewidth}p{0.13\linewidth}}
\toprule
Item & Outcome & Statements\\\midrule
A.1 & Surface existence and exact topology for every arrangement in scope; universal $1-n$ formula false & T1--T3, E1--E5, C.01--C.07\\
A.2 / L2 & Stable-state selector false; embedded diamond-and-band comparison proved & E6\\
A.3 / L3 & Proposed local image test false; complete image and graph reconstruction proved & P4--P6\\
A.4 & Disk sufficiency, injectivity, and cluster reversal proved, with exact degenerate qualification & P1--P3\\
L1 & Universal Euler and blow-up equivalence false; exact formula and equivalence criterion proved & E4--E5, C.04\\
L4 & Literal identification of tether and chord data false; marked two-tether correspondence proved & T4\\
\bottomrule
\end{tabular}
\end{center}

Every numbered statement has a status and a proof or an opened-source
derivation. Statement E1 uses a known theorem, while its equivariant
quotient consequence is established explicitly. The conditional
formulation E5 is applied only after T3 verifies every hypothesis.
The formulas give an effective determination: P4 enumerates the entire
finite extension set from the base chirotope, and hence supplies every
$c_j$. They do not posit that the answer depends on $n$ alone.

We distinguish the usual nonorientable genus $h=2-\chi$, the number of
projective-plane summands, from the manuscript parameter $g=1-\chi$.
Thus $h=g+1$. The finite test with $\chi=-45$ has $h=47$ and $g=46$.

The work is AI-produced by GPT-6 Astra (OpenAI, ChatGPT Pro mode).
No mathematical contribution of this run is attributed to a person,
and we claim no priority. Literature authors are cited for their
published results. The supplied AI studies and transcriptions are
witnesses to check, not authorities; manuscript references below are
cote/scan locators, not an unverified identification with handwritten
page numbers.
